\section{The Complete Scoped Assurance Protocol}
\label{sec:theory-protocol}
A public body must authorize a model for a specified public function, not
merely establish that its privacy audit passed. We therefore specify a
government governance profile over MRAP, with institutional obligations and
privacy evidence as separate, conjunctive conditions. Within that broader
decision, the principal statistical guarantee,
Theorem~\ref{thm:statistics}, controls erroneous ordinary privacy RELEASE
recommendations across covered selection and repeated opportunities.
It requires valid ceilings, not simply completed attacks. Soundness alone
would permit always refusing; the constructive routes below show when
informative ceilings and decisive comparisons are possible, while the
observation-relative countermodels identify conditions under which the
available information cannot support them.

MRAP supplies the complete normative lifecycle connecting these claims to
review, authorization and mediated service. A risk recommendation and an
institutional authorization remain distinct outputs: neither a signed
recommendation nor offline replay authorizes service.
We specify parties, immutable records, progression and refusal states,
evidence semantics and transaction boundaries before stating the guarantees.
Appendix~\ref{app:algorithms} gives pseudocode and
Appendix~\ref{app:proofs} the written proofs; the specification does not assert
that Python realizes every normative step.

\subsection{Parties, adversaries, and trusted boundaries}
The policy authority fixes obligations and tolerances; the population steward
approves protected units, priors, and scope; the owner and configuration
generator register artifacts and candidates. Evidence workers collect
observations, assessors replay them, and the optimizer applies the frozen
selection rule. An authorizing authority records independent governance
review, the portfolio registry commits the release, the gateway mediates
access, and monitors and the incident authority suspend or revoke it. Approved
keys, conflicts, separation of duties and permitted combinations of roles
belong to the trust profile; different filenames do not establish independence.

A \emph{privacy adversary} observes the registered export, auxiliary information
and allowed adaptive transcript to maximize inference gain. A separate
\emph{protocol adversary} may fabricate evidence, alter or replay messages,
withhold records, race concurrent proposals, or corrupt actors. Authentication,
collision-resistant bindings, sound inference/certificate procedures, a
faithful linearizable registry, trustworthy time/status checks and complete
per-request mediation are distinct assumptions. They do not imply that a
corrupted evidence authority measures truthfully or that the declared world
model is complete.

\subsection{Government governance profile and institutional gates}
\label{sec:gov-profile}
The responsible public body approves a profile $P_G$, bound into the existing
policy, trust and context attachments of instance $I$. It names the agency,
accountable service owner, competent authorizing office and independent
reviewers, including delegation limits and conflicts. It fixes the public
purpose, intended and prohibited uses, affected population, jurisdiction and
applicable legal/policy instruments. These are profile commitments, not new
JSON schemas or claims that the reference code already checks government
compliance. Applicability is decided by the institution's competent authorities;
the profile is not a universal statement of public law.

Each mandatory obligation has an owner, required evidence, authorized reviewer,
due stage, review interval, expiry and escalation route. The institution
registers mandate/delegation review; impact and rights review; appropriate
fairness, robustness, cybersecurity and procurement assessments; and human
oversight with specified powers to intervene. It also registers notice and
public-transparency arrangements, accessible explanation and redress routes,
incident handling, monitoring, renewal and retirement responsibilities.
Publication may have a reasoned, competent-authority-approved exemption where
the governing profile permits one; withholding a record does not erase its
internal review obligation. An assessor cannot silently delete a requirement
or invent an exemption after observing an unfavorable result.

Let $M_G^k(t)$ contain all mandatory institutional obligations due by stage
$k$ at time $t$, including earlier obligations whose evidence must remain
current and period-identified review obligations whose deadlines have arrived. For each
$j$, a bounded registered checker $V_j$ consumes scope-bound evidence $e_j$,
the instance, candidate, authoritative status $\Sigma$ and time $t$, returning
PASS, FAIL or UNKNOWN. Define
\[
 G^k(I,c,t)=\bigwedge_{j\in M_G^k(t)}
   [V_j(e_j;I,c,\Sigma,t)=\mathrm{PASS}].
\]
Evidence binds its issuer's competence/delegation, source material, rationale,
scope, conditions and current review status. For institutional judgments,
PASS records the prescribed competent attestation and supporting checks;
it is not an algorithmic proof that the judgment is legally or ethically
correct. Missing, expired, conflicting or unverifiable evidence is not PASS.
The named authority approves the mandatory roster and its applicability;
an empty caller-supplied roster cannot replace that profile.
Every successful progression edge checks $G^k$ for its destination before persistence,
using the authoritative time and status at that transition. Registration,
plan, evidence, assessment, selection, review, commitment and activation
therefore cannot skip a due institutional obligation. Later-due evidence
is not demanded as if its underlying activity had already occurred.
Refusal, denial, suspension and revocation retain their original authority
checks and do not require $G^k$ to pass: a failed obligation must not prevent
the protocol from stopping service.

Selection adds $G^{\rm select}$ to the existing feasible-set conditions.
Independent governance review is then obtained before commitment, rather
than assumed to exist at selection. At the authorization transaction define
\[
\begin{aligned}
 \operatorname{GovCommitOK}(I,c,t)={}&\operatorname{PrivEligible}(I,c)\\
 &{}\land G^{\rm commit}(I,c,t)\\
 &{}\land\operatorname{MRAPReady}(I,c,t),
\end{aligned}
\]
where PrivEligible is the privacy-disposition/mandatory-CLEAR rule below and
MRAPReady includes every other original selection, review and commit guard.
Only a successful atomic commit with this predicate issues a government
authorization. Privacy RELEASE or RELEASE-WITH-RISK alone never does.
Risk acceptance cannot waive a mandatory public-law or rights safeguard,
change the institutional roster, or substitute utility for required review.
An institutional refusal is recorded separately from the privacy interval:
a model can have acceptable evidenced privacy and still be unauthorized.

At activation and each service grant, require $G^{\rm serve}$, including
applicable intended-use, oversight, current-review and renewal conditions.
Before commit, future monitoring or redress obligations require approved,
resourced arrangements and responsible owners, not an assertion that all
future conduct has already occurred. Their due-time evidence is checked when
due; unknown status, a missed renewal or withdrawal of authority prevents new
grants. Changed purpose, authority, policy or material impact requires a fresh
instance and applicable reassessment; retirement ends service eligibility.
These gates enforce recorded institutional conditions under the trust model.
The six guarantees below do not become proofs of lawful government,
substantive fairness, effective redress or democratic legitimacy.

\subsection{Immutable context, messages, and evidence}
An instance $I$ fixes an identifier, expected registry head $h$, policy and
population snapshots, artifact and complete-interface commitments, finite
candidate set $\mathcal C$, nonempty in-scope sets $T_c$, utility requirements,
evidence plan, trust profile and expiry. Each game states access level,
auxiliary knowledge, query budget, adaptivity, prior, protected unit and metric.
The interface inventory covers weights, precision, preprocessing, outputs,
explanations, retrieval/tools/memory where applicable, metadata, logs and
related releases. An architecture label cannot replace this inventory.

The approved finite scenario universe $\mathcal D$ must satisfy
\[
 \mathcal D=\mathcal O\mathbin{\dot\cup}
                 \mathop{\dot\bigcup}_{t}\mathcal S_t,
\]
where exclusions $\mathcal O$ have reasons and every in-scope cell has an owner
and decision evidence. Channel-to-scenario witnesses include \emph{joint}
observations. Overlapping threat names may be converted to membership-profile
atoms of total predicates; the all-false residual atom stays visible.
This partitions supplied identifiers, not all real-world attacks.

Each message contains version/type, instance and message identifiers,
issuer/role, issue/expiry times, predecessor digests and payload. Its signature
covers a domain-separated digest of the canonical entire envelope. A consumer
captures source bytes once for hashing and parsing, checks type, role,
signature, freshness, predecessor state and exact instance binding, then
replays semantics. Signature validity alone is insufficient.
The evidence inventory binds data manifests, training configuration, code,
randomness records, completion metrics and final artifact hashes where policy
requires them. Pretrained artifacts retain their actual provenance limits.
Training start/finish are evidenced provenance events, not new MRAP phase
names; authentic records alone do not prove that training occurred.

For each $(c,t)$ use finite secret, observation and action spaces $S,Y,A$,
prior $\pi$, row-stochastic channel $K$ and gain $g:S\times A\to[0,1]$.
The complete joint transcript has value
\[
 \theta_{ct}=V_g(\pi,K)
   =\sum_y\max_a\sum_s\pi(s)K(y\mid s)g(s,a).
\]
This is posterior $g$-vulnerability \citep{alvim2012gain}. The definition and
its decision-theoretic interpretation are established foundations.
Infinite/continuous interfaces need a justified extension or an enforced
finite abstraction. For uniform binary guessing, success is $\theta$ and
normalized advantage is $2\theta-1$; these units must not be mixed with
attribution loss, canary exposure or another gain.

Admit evidence by \emph{direction}: justified lower endpoint, exact value,
valid analytic/statistical upper endpoint, or descriptive screen without
decision authority. For the same game and valid simultaneous coverage,
\[
 \begin{gathered}
 L=\max\{b,\ell_1,\ldots,\ell_r\},\quad U=\min\{1,u_1,\ldots,u_q\},\\
 b=\max_a\sum_s\pi(s)g(s,a).
 \end{gathered}
\]
Selection of these extrema must be covered. Missing optional endpoints may
use $b$ and one only after required evidence and control obligations are met.
A missing required experiment, unsupported endpoint, $L>U$ or invalid context
is a refusal, not a manufactured uncertainty interval. Gap is $U-L$; a
crossing is $L\le\tau<U$. Comparisons use exact represented rationals.
For incremental risk $V_g(\pi,K)-b$, the prior-only floor is zero and both
absolute endpoints must first have $b$ subtracted (then be bounded by the
metric's range). A prior upper cap is not an attained baseline. Ignoring $Y$
and choosing a prior-optimal action proves $V_g(\pi,K)\ge b$; consequently a
valid absolute ceiling below $b$ is inconsistent, not a reason to lower $L$.

\subsection{Integrated safeguards: repair or justified refusal}
\label{sec:admission}
The following O1--O7 rules implement the evidence conditions \emph{inside}
plan approval, admission, decision reduction and serving. Together with their
stated scientific and trust premises they provide a sufficient construction,
not a characterization of every possible sound protocol. The countermodels
show why particular weaker claims fail in specified observation classes;
they do not establish that this exact rule set is universally necessary or
minimal. The abstract total operation $\operatorname{Admit}(I,c,t,E)$
returns either $\operatorname{Valid}(L,U,r,\Gamma)$ or
$\operatorname{Refuse}(e,\rho)$, where $r$ is an approved bound route,
$\Gamma$ its explicit premises, $e$ a
reason and $\rho$ a remedy or scoped obstruction. These are mathematical
outputs, not new JSON schemas. Each approved route has a fixed, terminating
certificate/procedure checker with bounded input/resource limits; unrecognized
routes, failed checks and resource exhaustion refuse. Its soundness theorem
and observational trust premises are stated separately. Totality is a
decision over registered checkable routes, not an oracle deciding arbitrary
scientific truth. Unsupported inputs never receive a passing validity predicate.

\begin{enumerate}
\item \textbf{O1: direction and completeness.}
An upper endpoint must come from an approved analytic/structural theorem,
exact channel calculation, complete finite-class inference, or valid
simultaneous channel confidence set. A tested floor cannot fill that role.
An explicitly permitted worst-case $U=1$ is valid only after every required
evidence obligation is fulfilled; it is not evidence of low risk.
\item \textbf{O2: selection-wide accounting.}
Freeze the inference procedure, selectable family and stopping rules, and
record a valid simultaneous event and lifetime allocation before observation.
All influencing failed, abandoned and retried families remain charged.
Uncovered selection or a merely numerical allocation is refused.
\item \textbf{O3: scope and observation adequacy.}
Check the approved partition including residual/exclusions, all declared
channel witnesses and policy-required evidence of the observation boundary.
Refuse claims beyond that justified boundary. A partition alone does not
certify that a real deployment has no hidden interface.
\item \textbf{O4: joint evidence.}
Admit only a bound on the entire allowed joint transcript and current
portfolio, or a sound applicable joint composition theorem. Marginal
certificates alone cannot satisfy this obligation. Rebase and replay on a
changed head; absent required joint evidence is refusal.
\item \textbf{O5: metric-preserving transfer.}
Verify game, prior, access, decoder and gain identity, the safe direction of
transfer, and all numerical/structural premises. Charge an approximate
transfer's proven residual in its supported metric. Expected gain is not
silently substituted for worst-case posterior or another policy loss.
\item \textbf{O6: non-overridable decision semantics.}
Recompute the ordered disposition from every admitted in-scope cell.
Any refusal, contradictory interval, failed mandatory obligation or
demonstrated threshold violation blocks. Verify exact acceptance bindings
for every crossing. Acceptance cannot alter endpoints or substitute for
a strict mandatory CLEAR in selection or commit.
\item \textbf{O7: independent semantic and live binding.}
Require the registered source-observation, control, replay and
attestation/trust evidence, not just a signed assertion of its existence.
Before each service, require authoritative current status and faithful
deployment measurement/binding. Truth beyond the justified measurement
and trust model is an explicit assumption, not a signature theorem.
\end{enumerate}

If a required route fails or evidence is missing, Admit returns Refuse and
O6 makes the outward disposition BLOCK; applicable lifecycle rules then
abort or require redesign. If all obligations pass but only an allowed
$[L,1]$ is available, Admit returns Valid: the interval expresses real
bounded uncertainty and goes to the crossing/acceptance rules. A missing
mandatory battery, control or joint witness can never be repaired by
selecting $U=1$.

A no-go annotation identifies the requested claim, restricted observation
class, quantified safe/unsafe worlds and proof premises. It either explains
refusal of that claim or why an otherwise admissible interval remains
uninformative. It is never a clearance certificate, evidence waiver or
global infeasibility certificate. New observations or stronger enforcement
may escape its premises, but require a newly justified context/plan.
Thus the protocol addresses all seven inference hazards by constructive
admission rules or explicit refusal boundaries; it does not claim that every
deployment can supply the required evidence.

\subsection{State transitions and invalidation}
Table~\ref{tab:lifecycle} specifies the full normative machine. Every progression
consumes the immediately preceding authenticated records and rechecks their
continuing validity. Registration of an already trained model cannot
retrospectively preregister its training or audit selection. Plans and error
allocations precede the observations they purport to cover.

\begin{table*}[t]
\centering\small
\caption{Full normative MRAP lifecycle. Gates are conjunctive; records describe obligations, not implemented services.}
\label{tab:lifecycle}
\begin{tabular}{>{\raggedright\arraybackslash}p{0.26\textwidth}>{\raggedright\arraybackslash}p{0.23\textwidth}>{\raggedright\arraybackslash}p{0.43\textwidth}}
\toprule
Transition & Record and responsible role & Necessary condition beyond valid envelopes \\
\midrule
DRAFT $\to$ REGISTERED & Registration; owner & Purpose, parties, artifact, full interface, population, policy and current head bound (G0--G3). \\
REGISTERED $\to$ PLAN\_FROZEN & EvidencePlan; assessor and required approvers & Complete candidate/threat roster, controls, worker/source bindings, stopping/selection rule and prospective error allocation (G4--G5). \\
PLAN\_FROZEN $\to$ EVIDENCE\_FROZEN & EvidenceBundles; workers & Collection closed; required training provenance, failures/exclusions, controls and observed source bindings replayed (G6). \\
EVIDENCE\_FROZEN $\to$ ASSESSED & AssessmentReport; assessor & Every required cell evaluated; evidence direction, scope, arithmetic, transfer and joint portfolio checked (G7--G9). \\
ASSESSED $\to$ OPTIMIZED & OptimizationReport; optimizer & Mandatory privacy and utility constraints first; deterministic selection second; complete portfolio snapshot bound (G10). \\
OPTIMIZED $\to$ COMMIT\_PENDING & AuthorizationCommitRequest; authorizer & Independent review, objections, conflicts, conditions, expiry and incident/retirement ownership complete; no mandatory gate overridden (G11). \\
COMMIT\_PENDING $\to$ AUTHORIZED & AuthorizationReceipt; registry & Expected-head CAS, unused nonce, live evidence, complete portfolio and all budget deltas committed atomically (G12). \\
AUTHORIZED $\to$ ACTIVE & ActivationReceipt; gateway & Authoritative inclusion/current status, measured bytes/interface/controls and live lease verified (G13). \\
ACTIVE $\to$ SUSPENDED & Suspension/incident record; gateway/monitor & Material change, drift, budget/lease loss or unknown status stops new access (G14). \\
AUTHORIZED, ACTIVE, SUSPENDED $\to$ EXPIRED & Expiry record & Deadline elapsed; unused authorization and new service denied. \\
Nonterminal $\to$ REVOKED & Authorized revocation record & Registry records revocation; gateway denies new access. \\
Pre-authorization $\to$ REDESIGN\_REQUIRED & Remediation record & Missing/remediable evidence or incomplete feasible search; changed proposal needs a new instance. \\
ASSESSED/OPTIMIZED $\to$ REJECTED & Infeasibility/completeness certificate & No feasible candidate in a certified complete search; a searched subset is insufficient. \\
Pre-authorization $\to$ ABORTED & Failure record & Stale head, invalid/replayed/conflicting message, timeout or failed commit; no authorization issued. \\
\bottomrule
\end{tabular}
\end{table*}

REVOKED, EXPIRED, REJECTED and ABORTED are terminal. REDESIGN\_REQUIRED and
SUSPENDED supply no shortcut to service; a changed or reassessed proposal
uses a fresh instance. A same-ID, same-digest retry may return its recorded
receipt, but cannot extend a lease, undo revocation or bypass current service
checks. Same ID with different content is a violation. Any immutable change
invalidates downstream applicability; old signatures remain historical bytes,
not approvals of the new instance. Stale-head retries rebase and repeat
affected portfolio, budget, selection and review checks. Every evidence
family influencing later selection consumes its registered statistical
allocation, including failed or abandoned attempts; abort does not refund
an already observed experiment.

\paragraph{Persistent disclosure and service revocation}
Let $H_k$ conservatively contain every release potentially observed by the
declared adversary by step $k$, including revoked and expired releases, and
let $A_k$ be the active service roster. Require $A_k\subseteq H_k$ and
$H_k\subseteq H_{k+1}$. Portfolio evidence covers the joint observations in
$H_k$ plus the candidate, with adversary auxiliary knowledge and cross-context
links preserved. Suspension, revocation and retirement may shrink $A_k$, never
$H_k$. Before a new observable grant, atomically record its exposure or a
conservative reservation in $H$ with the budget/status transaction; if delivery
is uncertain, retain the reservation. A renamed purpose, instance or composition
domain cannot reset an adversary's knowledge. Independence or separation across
domains needs its own applicable argument. The monotonicity claim follows by
induction because exposure transitions add identifiers and all other transitions
leave $H$ unchanged; it assumes a complete initial roster and faithful mediation.
It is a specified invariant, not a new machine-checked result.

\subsection{Recommendation, selection, and authorization}
For one candidate let $B$ contain invalid/missing decision evidence for any
in-scope threat, missing mandatory obligations/controls, inconsistent bounds,
inadmissible clearance, stale context and invalid supplied acceptance.
Include every $L_t>\tau_t$. Let $X=\{t:L_t\le\tau_t<U_t\}$. Return, in order:
BLOCK if $B\ne\varnothing$; RELEASE if $X=\varnothing$;
RELEASE-WITH-RISK if valid acceptance covers every crossing;
INCONCLUSIVE if every crossing has a specific applicable resolving
measurement or bound; otherwise BLOCK.
A resolving plan identifies its estimand, procedure, allocation and
decision-relevant target, not a promise that every near-boundary case will
resolve. It binds the current instance, identifies obtainable data and an
available fresh error allocation, and justifies why the proposed precision or
bound can change this decision. A spent allocation, unavailable data, or a
generic suggestion to rerun the same test is insufficient. At $L=\tau$, a
plan needing a strictly positive separation must justify that separation or
be refused. The current outward implementation has no such plan verifier and
therefore refuses unaccepted crossings; scientific planning text is advisory.
Missing required evidence is BLOCK in this outward facade. Internal
scientific appraisal may retain INCONCLUSIVE/HOLD with redesign; neither
permits authorization.

Acceptance requires policy permission, an externally trusted acceptor,
current validity, and exact instance, interval, threshold and crossing-ceiling
bindings. It records risk accepted, not a smaller risk. The facade preserves
the original assessment's admission status: a below-threshold endpoint alone
cannot promote a non-CLEAR assessment to ordinary RELEASE. Risk acceptance
leaves every mandatory assessment status unchanged.

The optimizer admits a candidate only when its disposition is RELEASE or
RELEASE-WITH-RISK, every mandatory assessment cell is CLEAR, and all utility,
controls, transfer and portfolio obligations pass. RWR additionally requires
that every crossing is policy-permitted optional risk with exact valid
acceptance. An INCONCLUSIVE candidate cannot enter the feasible set.
It then applies the frozen least-information criterion and tie-break.
Neither utility nor RELEASE-WITH-RISK converts an unresolved mandatory cell
to CLEAR. Optional risk acceptance, where permitted, cannot waive a mandatory
gate. Thus a facade disposition alone is never the commit predicate.
An infeasible incomplete search yields redesign, not global rejection.
For a government instance, the feasible set additionally requires
$G^{\rm select}$; the authorizer obtains newly due independent institutional
reviews, and the registry requires GovCommitOK at commitment. A failed or
unknown institutional check refuses authorization even if the privacy
recommendation is RELEASE.

Within one transaction the registry checks authorizer, expected head, nonce,
live evidence, complete joint portfolio and all remaining budgets, then
commits authorization, budget deltas and new head together. No usable receipt
may escape an uncommitted transition. The gateway independently verifies
current inclusion/status and measured deployment identity at activation.
Each new access has a service linearization point that jointly validates the
current authorization epoch/status, reserves required budgets and grants the
request against an immutable measured artifact/interface snapshot. This point
also checks $G^{\rm serve}$ for government instances and is serialized with
revocation; stale epochs and unknown status deny access.
Revocation prevents grants linearized afterward, not necessarily completion
of earlier grants or recovery of delivered weights or outputs.

\subsection{Central guarantee and supporting results}
The central statistical inference is from valid, selection-wide ceilings to
controlled false ordinary privacy RELEASE, not from successful tests to
overall governmental legitimacy. Institutional gates remain additional
authorization conditions, not privacy confidence endpoints.
Proposition~\ref{prop:decision} supplies the deterministic implication used
by Theorem~\ref{thm:statistics}. The remaining results specify the scope of
that implication, routes for obtaining or transferring its endpoints, and
the extra enforcement premises needed to extend it to service.
Appendix~\ref{app:proofs} proves these conditional statements; none proves
the truth of an imported scientific assertion.

\begin{proposition}[Determinacy and endpoint soundness]
\label{prop:decision}
Fixed inputs yield exactly one disposition. On joint interval coverage,
RELEASE implies $\theta_t\le\tau_t$ for every in-scope threat; a blocking
floor implies $\theta_t>\tau_t$. Other BLOCK reasons establish inadmissibility,
not necessarily excessive risk.
\end{proposition}

\begin{theorem}[Selection- and lifetime-valid ordinary release]
\label{thm:statistics}
Before inspecting family $j$, allocate $\mathcal F_{j-1}$-measurable
$\alpha_j$, with $\sum_j\alpha_j\le\alpha$ almost surely. Its event $E_j$
covers every selectable candidate, threat, subgroup, portfolio and look,
with $\Pr(E_j^c\mid\mathcal F_{j-1})\le\alpha_j$. If all used ceilings are
sound deterministic bounds or belong to these families, and context,
collection and reduction premises hold, then
\[
 \Pr(\exists k,c,t:\ d_{ck}=\Release,\
                    \theta_{ctk}>\tau_{ctk})\le\alpha .
\]
\end{theorem}
Independence between families is unnecessary. This frequentist guarantee
excludes RELEASE-WITH-RISK and is not a posterior safety probability.
An allocated number without a valid procedure does not establish $E_j$.
Heterogeneous audits enter this argument through justified endpoints in their
own registered games. Their raw scores need not share a scale and are not
pooled: the union bound controls failures of the applicable ceiling claims,
while each risk is compared with its own tolerance. The result can therefore
combine different admissible evidence routes without treating a membership
score, canary exposure and provenance signal as interchangeable evidence.

\paragraph{Supporting scope, transfer, and enforcement results}
The next three results distinguish coverage of declared scenarios, validity
of transferred risk bounds, and preservation of the recommendation's context
through service. Their premises complement statistical coverage; none can
replace it or establish the adequacy of an unobserved deployment boundary.

\begin{proposition}[Relative scope completeness]
\label{prop:scope}
The disjoint-union check gives exactly one assignment or reasoned exclusion
per declared scenario. It establishes neither the correctness of exclusions
nor completeness of the real-world universe.
\end{proposition}

\begin{theorem}[Finite safe-direction transfer]
\label{thm:transfer}
For state-independent stochastic $Q$, $K_r=K_aQ$ implies
$V_g(\pi,K_r)\le V_g(\pi,K_a)$. If
$\max_s d_{\rm TV}(K_r(\cdot\mid s),(K_aQ)(\cdot\mid s))\le\eta$, then
$V_g(\pi,K_r)\le\min\{1,V_g(\pi,K_a)+\eta\}$.
\end{theorem}
This classical Blackwell/TV argument concerns bounded expected gain and the
complete joint channel. Restricted decoders require closure under composition
with $Q$, including resource costs. Posterior maxima and fixed-FPR operating
points need separate guarantees.

\begin{theorem}[Conditional lifecycle integrity and progress]
\label{thm:lifecycle}
With correct envelope/guard evaluation, binding and role assumptions, a faithful
linearizable registry, trusted time/status and atomic status/budget/grant
mediation serialized with revocation, every
served request has an unchanged instance, all ordered mandatory predecessors,
a committed live authorization and matching measured interface. If service
additionally requires ordinary RELEASE and Theorem~\ref{thm:statistics}'s
premises, a served registered-threshold violation has probability at most
$\alpha$. A feasible instance reaches ACTIVE along the finite path if every
required actor eventually responds successfully, checks terminate, the head
has no conflicting external change before commit and all deadlines remain live.
\end{theorem}
Conditional progress is not availability under withholding, head churn or
expiring evidence. Existing Lean nonvacuity does not prove this liveness claim.

\paragraph{A constructive route to nontrivial decisions}
The following finite-game result shows that empirical evidence can supply a
ceiling when the adversary class is complete and sampling assumptions hold.
It is deliberately stronger than testing a chosen attack battery: it bounds
the best allowed decoder, not merely the attacks that happened to succeed.

\begin{theorem}[Ceilings from a complete finite attack class]
\label{thm:finite-ceiling}
Fix all $M$ deterministic decoders of a finite game, $0<\alpha<1$, and
$n\ge1$ IID held-out pairs from its registered joint law $\pi(s)K(y\mid s)$,
with gains in $[0,1]$. Let
$e_n=\sqrt{\log(2M/\alpha)/(2n)}$ and
$h=\max_d\widehat V(d)$. Then
$[\max\{b,h-e_n\},\min\{1,h+e_n\}]$ covers the unrestricted game value with
probability at least $1-\alpha$. If $|\theta-\tau|\ge\delta>0$ and
$2e_n<\delta$, endpoints resolve the correct side on that event.
\end{theorem}
Completeness means all mappings $Y\to A$, so $M=|A|^{|Y|}$ and the strategy
attaining $b$ is included. For an incomplete class the ceiling bounds only
that class. A sufficient sample size is
$n>2\log(2M/\alpha)/\delta^2$. Literal enumeration is exponential in $|Y|$,
but the unrestricted empirical maximum separates by observation and can be
computed from counts; exponential enumeration is not a computational lower
bound. The certified transcript alphabet and sufficient labeled sampling
remain obligations. This applies Hoeffding's inequality
\citep{hoeffding1963bounded}; it is a
constructive route, not evidence of scalability to current LLM interfaces.
Exact channel evaluation or a valid simultaneous channel confidence set
with $U=\sup_{K'\in\mathcal K}V_g(\pi,K')$ are alternative conditional routes.

\subsection{Constructive sufficiency and observation-relative obstructions}
\label{sec:gaps}
O1--O7 constrain the normative protocol above, rather than only a subsequent
discussion of limitations. Appendix~\ref{app:gaps} pairs the seven inference
hazards with constructive repairs or explicit observation-relative
obstructions. These are not seven independent impossibility theorems:
selection and acceptance errors admit direct protocol repairs, whereas
nonidentifiability concerns the information available to a verifier.
Theorem~\ref{thm:observations} shows that safe and unsafe worlds with verifier
observations within TV distance $\Delta$ cannot have acceptance probabilities
differing by more than $\Delta$. It explains exactly why weak summaries,
marginal-only accounts, hidden channels and authentic-but-false statements
cannot support the excluded stronger claims within the constructed world
pairs. Richer observations or additional justified restrictions can escape
those premises; their adequacy must be evidenced through O3/O7. The argument
does not rule out a different sufficient protocol or prove a universal
impossibility of privacy assurance.

The constructive routes are explicit: complete finite inference or structural
ceilings (O1), simultaneous nonrefundable error accounting (O2), declared
coverage plus justified observation boundaries (O3), joint assessment (O4),
metric-correct transfer (O5), and non-override plus live mediation (O6/O7).
Section~\ref{sec:gap-study} reports finite tests of their witnesses and repairs.
No fixed sampling budget can guarantee uniformly reliable decisive answers
arbitrarily close to a threshold (Appendix~\ref{app:threshold}); this limits
promises of progress, not the correctness of a fail-closed refusal.
